\section{Which Federation, and Why}
\label{sec:ch05-which-federation-and-why}

\index{Composite Schemas specification|(}%
Chapter~\ref{ch:do-you-need-this} said that federation had outgrown the company
that named it, and that the GraphQL Specification Working Group had a
subcommittee producing a common specification for it. It then said this book
builds on Apollo Federation v2 anyway. Here is the reasoning, because a reader
who has just been walked through one directive vocabulary and is about to
spend the rest of a book on it is owed it.

\index{Composite Schemas specification!has no finished edition}%
The first reason is that there is nothing finished to build on. The GraphQL
Composite Schemas Spec publishes exactly one edition, labeled \enquote{Working
Draft}, and the project as a whole is marked \enquote{Prerelease}. There is no
numbered release and never has been~\autocite{compositeschemasspec}. That is not
a criticism of the work. It is a statement about what a book can responsibly
teach: a draft that may change is a fine thing to follow and a poor thing to
print several hundred pages against.

\index{Composite Schemas specification!is not a renaming}%
The second reason is that it is not Apollo Federation v2 with the serial
numbers filed off, so the choice is real rather than cosmetic. Two differences
are enough to make the point. Apollo's key is a \code{FieldSet}; the draft's is
a \code{FieldSelectionSet}, a different scalar with its own grammar. And where
Apollo puts a reserved root field on every subgraph, the draft marks an ordinary
field:

\begin{graphqlsdl}
directive @lookup on FIELD_DEFINITION
\end{graphqlsdl}

\index{entities@\code{\_entities}!absent from the draft}%
There is no \code{\_entities} in that document at all. A subgraph written to the
draft says \enquote{this existing field can be used to look an entity up}, where
a subgraph written to Apollo Federation says \enquote{here is a reserved field
that takes representations}. Those are different designs with different failure
modes, and half of part~III of this book is about the failure modes of the
second one.

\index{Fusion!cannot span both versions}%
The third reason is the one that actually decided it. The .NET implementation
of the draft is Fusion, and Fusion 16 is a rewrite rather than an upgrade.
Michael Staib, who founded ChilliCream, puts it plainly on their blog:
\enquote{When we set out to rewrite Fusion, we had three main goals: remove the
constraints imposed by Hot Chocolate, achieve top-tier performance with .NET,
and make the gateway feel like a natural extension of your ASP.NET Core app,
rather than a black box configured with YAML}~\autocite{staib2026fusion16}. The
same post says Fusion 16 implements the current version of the Composite Schema
Specification.

\index{HotChocolate.ApolloFederation@\code{HotChocolate.ApolloFederation}!spans three majors}%
A reader still on Hot Chocolate~14 cannot follow any of that, because the
Fusion that shipped in the 14 era is a different product that stopped receiving
releases and was replaced by differently named packages. \code{HotChocolate.ApolloFederation}
was not. It has an unbroken release line across all three major versions this
book could have covered, which makes Apollo Federation v2 the only federation
contract a reader on either of the two it does cover can write to. That is the
decision, and it is one about readers rather than about the merits of two
specifications.

\index{Composite Schemas specification!extends rather than replaces}%
None of which is a bet against the draft. The working group's own summary of its
August 2026 meeting records two things: that the specification's
final name will be the GraphQL Federation Specification, and that it
\enquote{would extend Apollo Federation rather than replace
it}~\autocite{compositeschemaswg202608}. I attach a caveat to that source rather
than to the claim, because the document carries one itself: it is an
automatically generated meeting summary and says so at the top. The rename had
not reached the specification when I checked on 23 August 2026: the published
document still called itself the Composite Schemas Spec.

\index{Apollo Federation!is the thing to learn either way}%
So the vocabulary in this chapter is worth learning on either outcome. If the
draft lands as an extension of what Apollo published, the model you are about to
read is the base it extends. If it diverges further, you will have run a
federated graph in production and will be in a position to judge the difference,
which is more than you can do from either specification's text.
\index{Composite Schemas specification|)}%
